\section{Sistematika Laporan}\label{sec:sistematika}

Laporan Tugas Akhir ini disusun dengan sistematika sebagai berikut:

\textbf{Bab I Pendahuluan} memaparkan konteks dan motivasi tugas akhir melalui latar belakang yang menjelaskan DecMed, keterbatasannya terhadap PGHD, dan tantangan integrasi yang dihadapi.
Bab ini juga memuat rumusan masalah, tujuan dan ukuran keberhasilan, batasan masalah, serta metodologi tugas akhir.

\textbf{Bab II Studi Literatur} membahas landasan teori dan studi terdahulu yang mendukung perancangan solusi. 
Subbab yang dibahas meliputi konsep rekam medis elektronik dan permasalahannya, \textit{Distributed Ledger Technology} dan karakteristik IOTA, arsitektur DecMed, konsep dan tantangan \textit{Patient-Generated Health Data} (PGHD), serta penelitian terkait yang menjadi pembeda tugas akhir ini.

\textbf{Bab III Analisis dan Perancangan Solusi} menjelaskan analisis permasalahan dan kebutuhan sistem, mencakup analisis kesenjangan DecMed terhadap PGHD, analisis alternatif arsitektur integrasi, analisis kebutuhan pipeline, kerahasiaan, integritas dan ketersediaan data, dan penjaminan \textit{data provenance}. 
Bab ini juga memuat rancangan solusi yang menjawab dan dirancang untuk menyelesaikan permasalahan yang sudah dianalisis sebelumnya.

\textbf{Bab IV Implementasi dan Pengujian} menjelaskan lingkungan pengembangan dan pengujian, detail implementasi prototipe sistem berdasarkan rancangan pada Bab III, serta hasil pengujian fungsional dan non-fungsional yang digunakan untuk mengevaluasi ketercapaian indikator keberhasilan yang telah ditetapkan.

\textbf{Bab V Kesimpulan dan Saran} merangkum hasil tugas akhir dalam bentuk kesimpulan yang menjawab setiap rumusan masalah, serta memberikan saran untuk pengembangan lebih lanjut berdasarkan keterbatasan yang ditemukan selama pengerjaan tugas akhir.